\documentclass[12pt,a4paper]{article}

\usepackage[T1]{fontenc}
\usepackage[utf8]{inputenc}
\usepackage{helvet}
\renewcommand{\familydefault}{\sfdefault}
\usepackage[spanish]{babel}
\usepackage{geometry}
\usepackage{booktabs}
\usepackage{array}
\usepackage{tabularx}
\usepackage{fancyhdr}
\usepackage{titlesec}
\usepackage{parskip}
\usepackage{microtype}
\usepackage{setspace}
\usepackage{xcolor}
\usepackage{lastpage}

\geometry{
  top=2.8cm,
  bottom=3cm,
  left=3cm,
  right=3cm,
  headheight=1.8cm
}

\definecolor{negro}{RGB}{0,0,0}
\definecolor{grismedio}{RGB}{100,100,100}
\definecolor{grisclaro}{RGB}{220,220,220}

\titleformat{\section}
  {\normalfont\large\bfseries\color{negro}}
  {}
  {0em}
  {}
\titlespacing{\section}{0pt}{20pt}{8pt}

\titleformat{\subsection}
  {\normalfont\normalsize\bfseries\color{negro}}
  {}
  {0em}
  {}
\titlespacing{\subsection}{0pt}{14pt}{4pt}

\pagestyle{fancy}
\fancyhf{}
\renewcommand{\headrulewidth}{0pt}
\fancyhead[L]{%
  \footnotesize\color{grismedio}%
  Project Charter — Plataforma Analítica de Mortalidad End to End%
}
\fancyhead[R]{%
  \footnotesize\color{grismedio}%
  Página \thepage\ de \pageref{LastPage}%
}
\fancyfoot[C]{%
  \color{grisclaro}\rule{\linewidth}{0.4pt}\\[3pt]
  \footnotesize\color{grismedio}CONFIDENCIAL%
}

\fancypagestyle{primerapagina}{%
  \fancyhf{}
  \renewcommand{\headrulewidth}{0pt}
}

\setlength{\parindent}{0pt}
\setlength{\parskip}{7pt}
\onehalfspacing

\begin{document}
\thispagestyle{primerapagina}

%% ─── PORTADA ─────────────────────────────────────────────────────────────────
\vspace*{0.4cm}
{\color{grisclaro}\rule{\linewidth}{1.2pt}}
\vspace{6pt}

\begin{center}
  {\fontsize{9}{11}\selectfont\color{grismedio}%
   \MakeUppercase{Universidad de San Carlos de Guatemala}\\[2pt]
   Facultad de Ingeniería \textendash\ Escuela de Ciencias y Sistemas}\\[10pt]
  {\fontsize{16}{20}\selectfont\bfseries\color{negro}
   Project Charter\\[4pt]
   Plataforma Analítica de Mortalidad End to End}
\end{center}

\vspace{6pt}
{\color{grisclaro}\rule{\linewidth}{1.2pt}}
\vspace{14pt}

{\small
\begin{tabular}{@{}p{3.8cm} p{9.8cm}@{}}
  \color{grismedio}Curso            & Seminario de Sistemas 2, Sección A \\[3pt]
  \color{grismedio}Catedrático      & Ingeniero Marlon Orellana \\[3pt]
  \color{grismedio}Tutor académico  & Estuardo Sebastián Valle Bances \\[3pt]
  \color{grismedio}Período          & Escuela de Vacaciones 2026 \\[3pt]
  \color{grismedio}Equipo consultor & Alberto Hernández, Fabio Hernández, Luis Morales \\[3pt]
  \color{grismedio}Versión          & 1.0 \\[3pt]
  \color{grismedio}Fecha            & 25 de junio de 2026 \\
\end{tabular}
}

%% ─── PROPÓSITO Y JUSTIFICACIÓN ───────────────────────────────────────────────
\newpage
\section{Propósito y Justificación}

El presente documento constituye el acta de constitución del proyecto, en adelante Project
Charter, correspondiente a la consultoría simulada bajo el modelo de Naciones Unidas y el
Programa de las Naciones Unidas para el Desarrollo, encargada de diseñar e implementar una
plataforma analítica de mortalidad de extremo a extremo para el Ministerio de Salud Pública y
Asistencia Social de Guatemala. El proyecto se justifica en la necesidad institucional de
comprender la evolución de los patrones de mortalidad entre el período previo y el período
posterior a la pandemia de COVID-19, con el fin de orientar decisiones de política pública y
asignación de recursos. Este documento formaliza el alcance, los objetivos, la estructura de
desglose del trabajo, el cronograma, la matriz de responsabilidades y el plan de gestión de
riesgos y calidad del proyecto en su totalidad.

%% ─── STAKEHOLDERS ───────────────────────────────────────────────────────────
\section{Partes Interesadas}

Las partes interesadas identificadas para este proyecto y su nivel de influencia sobre las
decisiones del mismo se detallan a continuación.

\vspace{6pt}
{\small
\begin{tabularx}{\linewidth}{@{} X l l X @{}}
  \toprule
  \textbf{Parte interesada} & \textbf{Rol} & \textbf{Influencia} & \textbf{Interés principal} \\
  \midrule
  Ministerio de Salud Pública y Asistencia Social de Guatemala
    & Cliente
    & Alta
    & Obtener evidencia analítica sobre mortalidad para orientar política pública \\[5pt]
  Programa de las Naciones Unidas para el Desarrollo
    & Patrocinador / Acompañante
    & Alta
    & Garantizar que la consultoría cumpla estándares internacionales de gobernanza
      y ética de datos \\[5pt]
  Ingeniero Marlon Orellana
    & Catedrático evaluador
    & Alta
    & Verificar el cumplimiento del alcance técnico y académico del encargo \\[5pt]
  Estuardo Sebastián Valle Bances
    & Tutor académico
    & Media
    & Orientar las decisiones técnicas del equipo y resolver ambigüedades en los
      términos de referencia \\[5pt]
  Equipo consultor
    & Ejecutor
    & Media
    & Entregar una plataforma de datos reproducible, auditable y bien documentada \\
  \bottomrule
\end{tabularx}
}

%% ─── OBJETIVOS ───────────────────────────────────────────────────────────────
\section{Objetivos del Proyecto}

El objetivo general consiste en diseñar, implementar y demostrar una plataforma de datos de
punta a punta que integre fuentes heterogéneas de mortalidad, las gobierne bajo una arquitectura
por capas hacia un repositorio analítico de tipo Data Warehouse, y entregue evidencia
comparativa entre los períodos previo y posterior a la pandemia que sustente recomendaciones de
política para la institución cliente.

De este objetivo general se derivan los siguientes objetivos específicos:

\vspace{4pt}
{\small
\begin{tabular}{@{} >{\bfseries}p{0.6cm} p{13cm} @{}}
  OE1 & Diseñar e implementar procesos de extracción, transformación y carga sobre fuentes de
        datos heterogéneas tanto locales como en la nube. \\[5pt]
  OE2 & Construir una arquitectura por capas con tablas de aterrizaje, de datos conformados y
        de modelo dimensional. \\[5pt]
  OE3 & Consolidar la información en un repositorio analítico tanto en la nube como en una
        instancia local, demostrando interoperabilidad entre ambas. \\[5pt]
  OE4 & Aplicar aprendizaje automático a un problema real de mortalidad. \\[5pt]
  OE5 & Producir visualización analítica en al menos dos herramientas de inteligencia de
        negocio distintas. \\[5pt]
  OE6 & Garantizar el manejo ético, anonimizado y agregado de los datos sensibles en
        cumplimiento de la normativa europea de protección de datos. \\[5pt]
  OE7 & Justificar técnicamente ante el cliente cada decisión de ingeniería adoptada. \\
\end{tabular}
}

%% ─── ALCANCE ─────────────────────────────────────────────────────────────────
\section{Alcance del Proyecto}

El alcance del proyecto comprende las tres fases acumulativas establecidas en los términos de
referencia del encargo. La primera fase abarca la identificación, estructuración, limpieza e
ingesta de los datos crudos hacia una zona de aterrizaje, con su correspondiente diccionario de
datos y plan de anonimización. La segunda fase abarca la transformación de dichos datos a través
de una arquitectura por capas hasta su consolidación en un modelo dimensional cargado
simultáneamente en un repositorio en la nube y en una réplica local interoperable. La tercera
fase abarca la aplicación de aprendizaje automático sobre el repositorio consolidado, la
producción de visualización analítica en múltiples herramientas de inteligencia de negocio, y la
entrega de un informe final de consultoría con recomendaciones de política pública basadas en
evidencia.

Queda fuera del alcance del proyecto la implementación de la plataforma en un entorno de
producción real para la institución cliente, dado que el encargo constituye una prueba de concepto
con fines exclusivamente educativos.

%% ─── SUPUESTOS Y RESTRICCIONES ───────────────────────────────────────────────
\section{Supuestos y Restricciones}

\subsection{Supuestos}

Los siguientes supuestos se declaran como verdaderos para efectos de la planificación y
ejecución del proyecto. Cualquier cambio en estas condiciones deberá evaluarse como un riesgo
activado y gestionarse mediante el plan de gestión de riesgos.

\vspace{4pt}
{\small
\begin{tabular}{@{} >{\bfseries}p{0.6cm} p{13cm} @{}}
  S1 & Las fuentes públicas de datos de mortalidad del Instituto Nacional de Estadística de
       Guatemala, la Organización Mundial de la Salud y el Banco Mundial permanecerán
       disponibles y accesibles durante todo el horizonte del proyecto. \\[5pt]
  S2 & La plataforma Databricks estará disponible en su modalidad de cómputo sin
       aprovisionamiento de infraestructura dedicada durante las fases de transformación y
       aprendizaje automático. \\[5pt]
  S3 & El equipo consultor contará con acceso continuo a los entornos de nube e instancia
       local requeridos para la arquitectura híbrida del repositorio analítico. \\[5pt]
  S4 & El diccionario de variables institucional proporcionado por la fuente original es
       suficiente para resolver las ambigüedades semánticas identificadas durante el análisis
       exploratorio. \\[5pt]
  S5 & Los términos de referencia del encargo autorizan explícitamente el uso de regresión
       logística y regresión de cresta como únicas técnicas de aprendizaje automático
       aplicables en la fase analítica. \\
\end{tabular}
}

\subsection{Restricciones}

\vspace{4pt}
{\small
\begin{tabular}{@{} >{\bfseries}p{0.6cm} p{13cm} @{}}
  R1 & El horizonte de ejecución del proyecto está fijado en cuatro meses calendario, sin
       posibilidad de extensión, dado que corresponde al período de la Escuela de Vacaciones
       2026 de la Facultad de Ingeniería. \\[5pt]
  R2 & El presupuesto total del proyecto está limitado a USD~20{,}278.21 bajo esquema de
       suma alzada, conforme a la propuesta financiera N.º PAM-2026-03. \\[5pt]
  R3 & El equipo consultor está conformado por exactamente tres integrantes, sin posibilidad
       de incorporar recursos adicionales durante la ejecución. \\[5pt]
  R4 & Los datos sensibles de salud deben manejarse en todo momento bajo los principios de
       anonimización conformes con el Reglamento (UE) 2023/2854, sin excepción. \\[5pt]
  R5 & La implementación en entorno de producción real queda fuera del alcance contratado,
       limitándose el entregable a una prueba de concepto reproducible y auditable. \\
\end{tabular}
}

%% ─── PRESUPUESTO ─────────────────────────────────────────────────────────────
\section{Presupuesto Resumido}

El detalle completo de tarifas por rol, asignación de persona-mes por fase y costos de
infraestructura se presenta en la Propuesta Financiera N.º PAM-2026-03 (Anexo Financiero),
la cual forma parte integral del contrato. El siguiente cuadro resume los montos aprobados por
categoría.

\vspace{6pt}
{\small
\begin{tabularx}{\linewidth}{@{} X r @{}}
  \toprule
  \textbf{Categoría} & \textbf{Monto aprobado (USD)} \\
  \midrule
  Costo de personal (honorarios del equipo consultor) & \$19{,}265.61 \\
  Costo no personal (infraestructura en la nube y licencias) & \$1{,}012.60 \\
  \midrule
  \textbf{Total aprobado} (suma alzada, todo incluido) & \textbf{\$20{,}278.21} \\
  \bottomrule
\end{tabularx}
}

\vspace{8pt}
El pago se realizará contra la entrega y aceptación satisfactoria de cada uno de los cinco hitos
definidos en el cronograma, conforme al detalle porcentual de la hoja \textit{«Cronograma de
Pagos»} del Anexo Financiero. No se realizarán pagos periódicos fijos independientes de la
entrega de productos.

%% ─── EDT ─────────────────────────────────────────────────────────────────────
\newpage
\section{Estructura de Desglose del Trabajo}

La estructura de desglose del trabajo organiza el proyecto en tres paquetes de trabajo principales,
correspondientes a cada una de las tres fases acumulativas del encargo, subdivididos a su vez en
los entregables específicos exigidos por los términos de referencia.

\subsection{Paquete de Trabajo 1: Fundamentos de Datos}

\vspace{4pt}
{\small
\begin{tabularx}{\linewidth}{@{} l X @{}}
  \toprule
  \textbf{Código} & \textbf{Entregable} \\
  \midrule
  1.1 & Identificación y catálogo de fuentes de datos públicas de mortalidad \\[3pt]
  1.2 & Ingesta de datos desde fuentes heterogéneas locales y en la nube \\[3pt]
  1.3 & Zona de aterrizaje sin transformación destructiva \\[3pt]
  1.4 & Diccionario de datos y plan de anonimización preliminar \\[3pt]
  1.5 & Modelo de datos producido en herramienta de modelado \\[3pt]
  1.6 & Arquitectura de despliegue de la solución \\[3pt]
  1.7 & Pipelines de ingesta versionados en repositorio de control de versiones \\[3pt]
  1.8 & Documentación desplegada mediante integración continua \\
  \bottomrule
\end{tabularx}
}

\subsection{Paquete de Trabajo 2: Transformación y Almacenamiento}

\vspace{4pt}
{\small
\begin{tabularx}{\linewidth}{@{} l X @{}}
  \toprule
  \textbf{Código} & \textbf{Entregable} \\
  \midrule
  2.1 & Capa de datos conformados con reglas de limpieza auditables \\[3pt]
  2.2 & Modelo dimensional de esquema estrella con tabla de hechos y dimensiones \\[3pt]
  2.3 & Orquestación de la transformación mediante flujo de trabajo de tareas dependientes \\[3pt]
  2.4 & Repositorio analítico híbrido con interoperabilidad entre nube e instancia local \\[3pt]
  2.5 & Sistema de respaldo automatizado del repositorio en la nube hacia la instancia local \\[3pt]
  2.6 & Modelo de datos de hechos y dimensiones con su diagrama entidad relación \\
  \bottomrule
\end{tabularx}
}

\subsection{Paquete de Trabajo 3: Aprendizaje Automático y Business Intelligence}

\vspace{4pt}
{\small
\begin{tabularx}{\linewidth}{@{} l X @{}}
  \toprule
  \textbf{Código} & \textbf{Entregable} \\
  \midrule
  3.1 & Modelo de aprendizaje automático entrenado y evaluado sobre el repositorio analítico \\[3pt]
  3.2 & Visualización analítica en al menos dos herramientas de inteligencia de negocio \\[3pt]
  3.3 & Análisis comparativo entre el período previo y posterior a la pandemia \\[3pt]
  3.4 & Informe final de consultoría con recomendaciones de política pública \\[3pt]
  3.5 & Gestión del proyecto: acta de constitución, cronograma, estructura de desglose
        y matriz de responsabilidades \\[3pt]
  3.6 & Propuesta financiera de la consultoría \\[3pt]
  3.7 & Defensa oral con demostración en vivo de extremo a extremo \\
  \bottomrule
\end{tabularx}
}

%% ─── CRITERIOS DE ACEPTACIÓN ─────────────────────────────────────────────────
\newpage
\section{Criterios de Aceptación}

Para que cada fase sea considerada formalmente aceptada por la institución cliente, deberán
cumplirse los siguientes criterios, verificables de manera objetiva en el momento de la entrega.

\vspace{6pt}
{\small
\begin{tabularx}{\linewidth}{@{} l X @{}}
  \toprule
  \textbf{Fase / Entregable} & \textbf{Criterio de aceptación} \\
  \midrule
  Fase 1 — Zona de aterrizaje
    & Los datos crudos de todas las fuentes identificadas están ingestados sin
      transformación destructiva, el diccionario de variables está completo y el plan de
      anonimización preliminar ha sido revisado y aprobado por el tutor académico. \\[5pt]
  Fase 2 — Modelo dimensional
    & El repositorio analítico presenta cero registros huérfanos en la verificación de
      integridad referencial entre la tabla de hechos y cada dimensión, y la réplica local
      es consultable de manera independiente de la instancia en la nube. \\[5pt]
  Fase 3 — Aprendizaje automático
    & El modelo de clasificación alcanza un área bajo la curva ROC superior a 0.50 sobre
      el conjunto de prueba, y el modelo de pronóstico produce predicciones evaluables
      sobre el período 2023--2024 con métricas documentadas. \\[5pt]
  Fase 3 — Visualización analítica
    & Las visualizaciones están publicadas y son accesibles en al menos dos herramientas
      de inteligencia de negocio distintas, con conexión directa al repositorio analítico. \\[5pt]
  Informe final
    & El informe documenta la metodología completa, los hallazgos del análisis comparativo,
      las recomendaciones de política pública y las limitaciones del estudio, en formato
      entregable aprobado por el catedrático. \\[5pt]
  Defensa oral
    & La demostración en vivo cubre el recorrido completo de extremo a extremo desde la
      fuente original hasta el resultado analítico, sin interrupciones técnicas no
      justificadas. \\
  \bottomrule
\end{tabularx}
}

%% ─── CRONOGRAMA ──────────────────────────────────────────────────────────────
\newpage
\section{Cronograma General del Proyecto}

El proyecto se planificó bajo un horizonte de cuatro meses calendario, distribuido en tres fases
secuenciales y acumulativas, cada una con sus correspondientes hitos de entrega y defensa ante
el cliente simulado. Cada fase incluye una semana de margen interno para revisión de calidad
previa a su entrega formal.

\vspace{6pt}
{\small
\begin{tabularx}{\linewidth}{@{} X c c c c @{}}
  \toprule
  \textbf{Paquete de trabajo} & \textbf{Mes 1} & \textbf{Mes 2} & \textbf{Mes 3} & \textbf{Mes 4} \\
  \midrule
  1. Fundamentos de Datos
    & \multicolumn{2}{c}{Ejecución y cierre} & & \\[4pt]
  2. Transformación y Almacenamiento
    & & Inicio & \multicolumn{2}{c}{Ejecución y cierre} \\[4pt]
  3. Aprendizaje Automático y Business Intelligence
    & & & Inicio -- Ejecución & Cierre y defensa \\[4pt]
  Gestión del proyecto y aseguramiento de calidad
    & Continuo & Continuo & Continuo & Continuo \\
  \bottomrule
\end{tabularx}
}

\vspace{14pt}
\subsection{Hitos Principales}

\vspace{4pt}
{\small
\begin{tabularx}{\linewidth}{@{} l X c @{}}
  \toprule
  \textbf{Hito} & \textbf{Entregable asociado} & \textbf{Mes} \\
  \midrule
  Hito 1 & Entrega y defensa de la Fase 1: Fundamentos de Datos & 1 \\[3pt]
  Hito 2 & Entrega y defensa de la Fase 2: Transformación y Almacenamiento & 2 \\[3pt]
  Hito 3 & Tarea complementaria de integraciones con herramientas de procesamiento & 2 \\[3pt]
  Hito 4 & Entrega y defensa de la Fase 3: Aprendizaje Automático y Business Intelligence & 4 \\[3pt]
  Hito 5 & Entrega del informe final de consultoría y cierre del proyecto & 4 \\
  \bottomrule
\end{tabularx}
}

%% ─── MATRIZ RACI ─────────────────────────────────────────────────────────────
\newpage
\section{Matriz de Responsabilidades}

La matriz de responsabilidades se elaboró bajo la metodología RACI, que clasifica la
participación de cada integrante del equipo consultor en cuatro niveles: \textbf{R}~responsable
de la ejecución, \textbf{A}~autoridad que aprueba, \textbf{C}~consultado para insumos técnicos,
e \textbf{I}~informado sobre el avance.

\vspace{6pt}
{\small
\begin{tabularx}{\linewidth}{@{} X c c c @{}}
  \toprule
  \textbf{Actividad}
    & \textbf{Alberto H.}
    & \textbf{Fabio H.}
    & \textbf{Luis M.} \\
  \midrule
  Ingesta y arquitectura de datos               & R, A & C    & I    \\[3pt]
  Análisis exploratorio y reglas de transformación & R, A & C  & I    \\[3pt]
  Sistema de respaldo y auditoría               & R, A & I    & I    \\[3pt]
  Transformación ETL y modelo dimensional       & I    & R, A & C    \\[3pt]
  Orquestación del flujo de trabajo             & C    & R, A & I    \\[3pt]
  Modelado dimensional en herramienta de modelado & I  & R, A & C    \\[3pt]
  Aprendizaje automático                        & R, A & C    & C    \\[3pt]
  Visualización en herramientas de inteligencia de negocio & R & R & A \\[3pt]
  Documentación y gobernanza de datos           & C    & I    & R, A \\[3pt]
  Gestión del proyecto y aseguramiento de calidad & C  & C    & R, A \\
  \bottomrule
\end{tabularx}
}

%% ─── GESTIÓN DE RIESGOS ──────────────────────────────────────────────────────
\newpage
\section{Gestión de Riesgos}

Se identificaron los riesgos más relevantes observados durante la ejecución del proyecto, junto
con su probabilidad de ocurrencia, impacto sobre el cronograma o la calidad del entregable, y la
estrategia de mitigación adoptada en cada caso.

\vspace{6pt}
{\small
\begin{tabularx}{\linewidth}{@{} X c c X @{}}
  \toprule
  \textbf{Riesgo} & \textbf{Prob.} & \textbf{Impacto} & \textbf{Estrategia de mitigación} \\
  \midrule
  Bloqueo de acceso a fuentes públicas mediante mecanismos de protección automatizada
    & Alta & Medio
    & Automatización del acceso mediante herramientas de navegación programática que
      replican la interacción humana \\[6pt]
  Errores de interpretación semántica en variables de la fuente original
    & Media & Alto
    & Análisis exploratorio exhaustivo previo a la definición de reglas de transformación,
      con verificación cruzada contra el diccionario de variables institucional \\[6pt]
  Defectos de integridad en el modelo dimensional no detectados antes de la fase de
  explotación analítica
    & Media & Alto
    & Verificación sistemática de integridad referencial y de unicidad de llaves primarias
      en cada dimensión antes de avanzar a la fase siguiente \\[6pt]
  Limitaciones de las funcionalidades en vista previa pública de la plataforma de
  cómputo distribuido
    & Media & Alto
    & Identificación temprana de restricciones documentadas de la plataforma y adopción
      de bibliotecas de respaldo no sujetas a dichas restricciones \\[6pt]
  Ambigüedad en los términos de referencia respecto a la herramienta de aprendizaje
  automático exigida
    & Baja & Medio
    & Consulta directa y oportuna al tutor académico para resolver la ambigüedad antes
      de invertir esfuerzo en una implementación incorrecta \\[6pt]
  Concentración de tareas críticas en fechas próximas a la entrega final
    & Alta & Alto
    & Distribución explícita de responsabilidades entre los tres integrantes del equipo
      y priorización de entregables según su peso en la evaluación \\
  \bottomrule
\end{tabularx}
}

%% ─── PLAN DE CALIDAD ─────────────────────────────────────────────────────────
\newpage
\section{Plan de Aseguramiento de Calidad}

El plan de aseguramiento de calidad del proyecto se sustentó en cuatro mecanismos aplicados de
manera transversal a lo largo de las tres fases.

\vspace{4pt}
{\small
\begin{tabular}{@{} >{\bfseries}p{0.6cm} p{13cm} @{}}
  M1 & Toda regla de transformación aplicada sobre los datos se documentó con un código de
       referencia trazable hacia el hallazgo del análisis exploratorio que la originó, permitiendo
       auditar la justificación de cada decisión de limpieza. \\[6pt]
  M2 & Cada flujo de transformación incorporó una verificación automática de integridad
       referencial entre la tabla de hechos y sus dimensiones, con el criterio de aceptación de
       cero registros huérfanos en cada relación. \\[6pt]
  M3 & Todo cambio de código se validó de manera incremental antes de su incorporación al
       flujo de producción, mediante pruebas con datos sintéticos que replican la estructura real
       de los datos del proyecto, previo a su ejecución sobre el volumen completo de registros. \\[6pt]
  M4 & Cada ejecución del flujo de trabajo completo, así como cada entrenamiento de modelo de
       aprendizaje automático, quedó registrada en una tabla de control de auditoría con marcas
       de tiempo, volumen de registros procesados y métricas resultantes, permitiendo la
       trazabilidad completa de cualquier resultado presentado en el informe final. \\
\end{tabular}
}

%% ─── AUTORIZACIÓN ────────────────────────────────────────────────────────────
\newpage
\section{Autorización del Project Charter}

La firma del presente documento por las partes indicadas formaliza la aprobación del alcance,
los objetivos, el presupuesto y las condiciones de ejecución aquí establecidos, y autoriza al
equipo consultor a dar inicio a la primera fase del proyecto.

\vspace{32pt}

\noindent
\begin{tabular}{@{} p{7cm} @{\hspace{1.5cm}} p{7cm} @{}}

  \color{negro}\rule{6.5cm}{0.7pt} &
  \color{negro}\rule{6.5cm}{0.7pt} \\[8pt]

  \textbf{Alberto Hernández} &
  \textbf{Por la institución cliente} \\[2pt]
  \small Encargado del Proyecto &
  \small Ministerio de Salud Pública \\
  \small Equipo Consultor &
  \small y Asistencia Social de Guatemala \\
  &
  \small Programa de las Naciones Unidas \\
  &
  \small para el Desarrollo \\[24pt]

  \small Lugar y fecha: \dotfill &
  \small Lugar y fecha: \dotfill \\

\end{tabular}

%% ─── REFERENCIAS ─────────────────────────────────────────────────────────────
\newpage
\section{Referencias}

\begingroup
\setlength{\parskip}{6pt}
\hangindent=1.5em
Project Management Institute. (2021). \textit{A guide to the project management body of
knowledge (PMBOK guide)} (7th ed.). Project Management Institute.
\endgroup

\end{document}